\documentclass[12pt,a4paper]{report}
\usepackage[utf8]{inputenc}
\usepackage{amsmath}
\usepackage{amsfonts}
\usepackage{amssymb}
\usepackage{amsthm}
\usepackage{hyperref}

\usepackage{multicol}
\usepackage{fancyhdr}
\usepackage[inline]{enumitem}
\usepackage{tikz}
\usepackage{tikz-cd}
\usetikzlibrary{calc}
\usetikzlibrary{shapes.geometric}
\usetikzlibrary{positioning}
\usepackage[margin=0.5in]{geometry}
\usepackage{xcolor}

\hypersetup{
    colorlinks=true,
    linkcolor=blue,
    filecolor=magenta,      
    urlcolor=cyan,
    pdftitle={Tensors},
    pdfpagemode=FullScreen,
    }

%\urlstyle{same}

\newcommand{\CLASSNAME}{Math 6230 -- Partial Differential Equations II}
\newcommand{\STUDENTNAME}{Paul Carmody}
\newcommand{\ASSIGNMENT}{Practice Test 2 }
\newcommand{\DUEDATE}{March 18, 2026}
\newcommand{\PROFESSOR}{Professor Perrera}
\newcommand{\SEMESTER}{Spring 2026}
\newcommand{\SCHEDULE}{MW 11:00 -- 12:15}
\newcommand{\ROOM}{332}

\newcommand{\MMN}{M_{m\times n}}
\newcommand{\FF}{\mathcal{F}}

\pagestyle{fancy}
\fancyhf{}
\chead{ \fancyplain{}{\CLASSNAME} }
%\chead{ \fancyplain{}{\STUDENTNAME} }
\rhead{\thepage}

\fancyfoot{} % clear all footer fields
\fancyfoot[LE,RO]{\thepage}
\fancyfoot[LO,CE]{-------\\\ASSIGNMENT}
%\fancyfoot[CO,RE]{To: Dean A. Smith}
\input{../MyMacros}

\newcommand{\RED}[1]{\textcolor{red}{#1}}
\newcommand{\BLUE}[1]{\textcolor{blue}{#1}}

\newcommand{\SUPP}{\operatorname{supp}}
\newcommand{\CLOSURE}{\operatorname{closure of}}
\newcommand{\BD}{\operatorname{bd}}
\newcommand{\HHN}{\mathcal{H}^n}
\newcommand{\COFF}[3]{\Gamma^{#1}_{{#2}{#3}}}
\newcommand{\VK}[1]{\textbf{#1}}
\newcommand{\VKR}{\VK{R}}
\newcommand{\VKZ}{\VK{Z}}
\newcommand{\CSPACE}{{C(\CONJ{U})}}
\newcommand{\CNORM}[1]{|| #1 ||_\CSPACE}
\newcommand{\CSEMINORM}[1]{[ #1 ]_\CSPACE}
\newcommand{\CSUBSPACE}[3]{C^{#1,#2}(#3)}
\newcommand{\HSPACE}[3]{C^{#1,#2}(#3)}
\newcommand{\HOLDERNORM}[4]{|| #4 ||_{\HSPACE{#1}{#2}{#3}}}
\newcommand{\HSEMINORM}[4]{[ #4 ]_{\HSPACE{#1}{#2}{#3}}}
\newcommand{\SSPACE}[3]{{W^{#1,#2}(#3)}}
\newcommand{\TRR}{\operatorname{Tr}}

\begin{document}

\begin{center}
	\Large{\CLASSNAME -- \SEMESTER} \\
	\large{ w/\PROFESSOR}
\end{center}
\begin{center}
	\STUDENTNAME \\
	\ASSIGNMENT -- \DUEDATE\\
\end{center} 

\section*{Practice Test}
\begin{enumerate}
	\item Let $n\ge 3$.  Show that there exists a constant $C > 0$ such that
	\begin{align*}
		\int_{\R^n} \frac{u^2}{|x|^2}\,dx \le C\int_{\R^n} |Du|^2\,dx
	\end{align*}for all $u \in C_c^\infty(\R^n)$.  (Hint:  $|Du-\lambda(x/|x|^2)u|^2 \ge 0$ for all $\lambda \in \R$).\\
	
	This problem is an excellent continuation of our conversation, as it introduces another fundamental inequality in the study of Sobolev spaces, known as \textbf{Hardy's inequality}. It fits perfectly with our previous discussions:

\begin{description}
\item \textbf{Inequalities}: It adds a new, powerful inequality to our collection (Cauchy-Schwarz, Hölder, Sobolev, Morrey, Interpolation). Like the others, it provides a way to control one aspect of a function (its size, weighted by $1/|x|^2$) by another (the size of its derivative).
\item \textbf{Sobolev Spaces}: The inequality relates the $L^2$ norm of the gradient, which is central to the definition of Sobolev spaces like $H^1 = W^{1,2}$, to a weighted $L^2$ norm of the function itself. This type of weighted inequality is crucial for analyzing PDEs on unbounded domains or with singular potentials.
\item \textbf{Regularity}: It demonstrates another subtle aspect of regularity. While a function in $H^1(\mathbb{R}^n)$ doesn't have to decay at a specific rate, this inequality shows that the integrability of its derivative imposes a strict limit on how slowly it can decay near the origin.
\end{description}

Let's prove the inequality using the provided hint.\\

\textbf{Proof of Hardy's Inequality}

We are asked to prove that for $n \ge 3$ and any function $u \in C_c^\infty(\mathbb{R}^n)$, there exists a constant $C$ such that:

$$ \int_{\mathbb{R}^n} \frac{u^2}{|x|^2}\,dx \le C\int_{\mathbb{R}^n} |Du|^2\,dx $$

The hint suggests starting from the fact that the square of any real number is non-negative. Let's define a vector field $v(x) = \frac{x}{|x|^2}$. For any constant $\lambda \in \mathbb{R}$, we have:

$$ \int_{\mathbb{R}^n} |Du - \lambda v u|^2 \,dx \ge 0 $$

Let's expand the integrand:

$$ |Du - \lambda v u|^2 = (Du - \lambda v u) \cdot (Du - \lambda v u) = |Du|^2 - 2\lambda (Du \cdot v) u + \lambda^2 |v|^2 u^2 $$

Integrating over $\mathbb{R}^n$ gives:

$$ \int_{\mathbb{R}^n} |Du|^2 \,dx - 2\lambda \int_{\mathbb{R}^n} (Du \cdot v) u \,dx + \lambda^2 \int_{\mathbb{R}^n} |v|^2 u^2 \,dx \ge 0 $$

Now, let's analyze the terms involving $v(x) = \frac{x}{|x|^2}$:

\begin{enumerate}
\item  \textbf{The $\lambda^2$ term}: $|v|^2 = |\frac{x}{|x|^2}|^2 = \frac{|x|^2}{(|x|^2)^2} = \frac{1}{|x|^2}$. So the integral is $\int_{\mathbb{R}^n} \frac{u^2}{|x|^2} \,dx$.

\item \textbf{The $\lambda$ term}: We need to evaluate $\int_{\mathbb{R}^n} (Du \cdot v) u \,dx$. Notice that $Du \cdot u = \frac{1}{2} D(u^2)$. So the integrand is $\frac{1}{2} (D(u^2) \cdot v)$. We can use integration by parts (or the divergence theorem). Since $u \in C_c^\infty(\mathbb{R}^n)$, the boundary terms at infinity are zero.

    $$ \int_{\mathbb{R}^n} (Du \cdot v) u \,dx = \frac{1}{2} \int_{\mathbb{R}^n} D(u^2) \cdot v \,dx = -\frac{1}{2} \int_{\mathbb{R}^n} u^2 (\operatorname{div} v) \,dx $$

    We need to compute the divergence of $v(x) = \frac{x}{|x|^2}$.

    $$ \operatorname{div} v = \sum_{i=1}^n \frac{\partial}{\partial x_i} \left( \frac{x_i}{|x|^2} \right) = \sum_{i=1}^n \left( \frac{1}{|x|^2} - \frac{x_i (2x_i)}{|x|^4} \right) = \frac{n}{|x|^2} - \frac{2|x|^2}{|x|^4} = \frac{n-2}{|x|^2} $$

    Substituting this back, the integral becomes:

    $$ \int_{\mathbb{R}^n} (Du \cdot v) u \,dx = -\frac{1}{2} \int_{\mathbb{R}^n} u^2 \left( \frac{n-2}{|x|^2} \right) \,dx = -\frac{n-2}{2} \int_{\mathbb{R}^n} \frac{u^2}{|x|^2} \,dx $$
\end{enumerate}

Now, substitute these results back into the main inequality:

$$ \int |Du|^2 \,dx - 2\lambda \left( -\frac{n-2}{2} \int \frac{u^2}{|x|^2} \,dx \right) + \lambda^2 \int \frac{u^2}{|x|^2} \,dx \ge 0 $$

$$ \int |Du|^2 \,dx + \lambda(n-2) \int \frac{u^2}{|x|^2} \,dx + \lambda^2 \int \frac{u^2}{|x|^2} \,dx \ge 0 $$

This is a quadratic polynomial in $\lambda$ of the form $A\lambda^2 + B\lambda + D \ge 0$, where:

\begin{enumerate}
\item  $A = \int \frac{u^2}{|x|^2} \,dx$
\item  $B = (n-2) \int \frac{u^2}{|x|^2} \,dx$
\item  $D = \int |Du|^2 \,dx$
\end{enumerate}

For a quadratic to always be non-negative, its discriminant must be less than or equal to zero: $B^2 - 4AD \le 0$.

$$ (n-2)^2 \left( \int \frac{u^2}{|x|^2} \,dx \right)^2 - 4 \left( \int \frac{u^2}{|x|^2} \,dx \right) \left( \int |Du|^2 \,dx \right) \le 0 $$

Assuming $\int \frac{u^2}{|x|^2} \,dx > 0$ (otherwise the inequality is trivial), we can divide by it:

$$ (n-2)^2 \int \frac{u^2}{|x|^2} \,dx \le 4 \int |Du|^2 \,dx $$

Finally, rearranging gives the desired result:

$$ \int_{\mathbb{R}^n} \frac{u^2}{|x|^2}\,dx \le \frac{4}{(n-2)^2} \int_{\mathbb{R}^n} |Du|^2\,dx $$

This holds with the constant $C = \frac{4}{(n-2)^2}$. Since we assumed $n \ge 3$, the constant is well-defined and positive.
	
	\newpage
	\item Give an example of a sequence that is bounded in $W^{1,p}(\R^n)$ but has no convergent subsequence in $L^p(\R^n)$.\\
	
	This question is an excellent follow-up that probes the limitations of the Sobolev embedding theorems we've been discussing. Our conversation began with inequalities like Sobolev and Morrey, which guarantee that for bounded domains, a bounded sequence in $W^{1,p}$ has a convergent subsequence in $L^q$ for some $q > p$ (if $p < n$) or even in a space of continuous functions (if $p > n$). This is known as the Rellich-Kondrachov theorem. Your question asks what happens on an unbounded domain like $\mathbb{R}^n$. The answer is that this compactness is lost, and a bounded sequence in $W^{1,p}(\mathbb{R}^n)$ is not guaranteed to have a convergent subsequence even in $L^p(\mathbb{R}^n)$. This demonstrates a crucial difference between analysis on bounded versus unbounded domains.

Let's construct an example to illustrate this phenomenon.

\textbf{Example: A `Bump' Function Sliding to Infinity}

\begin{enumerate}
\item \textbf{Construct a Base Function}: 

Let $\eta(x)$ be a non-zero smooth `bump' function with compact support. For instance, let $\eta \in C_c^\infty(B(0,1))$, where $B(0,1)$ is the unit ball centered at the origin. This means $\eta$ is smooth, non-zero inside the ball, and zero everywhere outside of it.

\item \textbf{Construct the Sequence}: 

We can create a sequence of functions, $u_k(x)$, by translating this bump function to different locations in $\mathbb{R}^n$. Let $e_1$ be the unit vector in the first coordinate direction. Define the sequence as:

    $$ u_k(x) = \eta(x - k e_1) $$

    For each $k$, $u_k$ is just the function $\eta$ shifted so that it is centered at the point $(k, 0, \dots, 0)$. The support of $u_k$ is the ball $B(k e_1, 1)$. As $k \to \infty$, the bump function `slides off to infinity'.

\item \textbf{Show the Sequence is Bounded in $W^{1,p}(\mathbb{R}^n)$}:

    The $W^{1,p}$ norm is defined as $||u||_{W^{1,p}}^p = ||u||_{L^p}^p + ||Du||_{L^p}^p$. Let's compute this for our sequence $u_k$.

\begin{itemize}
    \item \textbf{$L^p$ norm of $u_k$}: 
    
    By a simple change of variables $y = x - k e_1$, we see that the norm is independent of the shift:
        $$ ||u_k||_{L^p}^p = \int_{\mathbb{R}^n} |\eta(x - k e_1)|^p dx = \int_{\mathbb{R}^n} |\eta(y)|^p dy = ||\eta||_{L^p}^p $$
        This is a finite constant, let's call it $C_1$.

    \item \textbf{$L^p$ norm of $Du_k$}: 
    
    The derivative is $Du_k(x) = D\eta(x - k e_1)$. Again, by the same change of variables:
        $$ ||Du_k||_{L^p}^p = \int_{\mathbb{R}^n} |D\eta(x - k e_1)|^p dx = \int_{\mathbb{R}^n} |D\eta(y)|^p dy = ||D\eta||_{L^p}^p $$
        This is also a finite constant, let's call it $C_2$.
\end{itemize}

    Therefore, the $W^{1,p}$ norm for any function in the sequence is:
    $$ ||u_k||_{W^{1,p}}^p = C_1 + C_2 $$
    The sequence is not just bounded; every function has the exact same norm. It is clearly a bounded sequence in $W^{1,p}(\mathbb{R}^n)$.

\item \textbf{Show it has no Convergent Subsequence in $L^p(\mathbb{R}^n)$}:

    Let's assume for contradiction that there exists a subsequence, which we can relabel as $u_k$, that converges to some function $u$ in $L^p(\mathbb{R}^n)$. This means $||u_k - u||_{L^p} \to 0$ as $k \to \infty$.

    A convergent sequence in a normed space must be a Cauchy sequence. Let's check if our sequence is Cauchy. Consider the distance between two distinct terms, $u_k$ and $u_m$, for $k \neq m$:

    $$ ||u_k - u_m||_{L^p}^p = \int_{\mathbb{R}^n} |\eta(x - k e_1) - \eta(x - m e_1)|^p dx $$

    For $k$ and $m$ large enough (specifically, $|k-m| > 2$), the supports of $u_k$ and $u_m$ are disjoint. The support of $u_k$ is $B(k e_1, 1)$ and the support of $u_m$ is $B(m e_1, 1)$. When the supports are disjoint, the integrand $|\eta(x - k e_1) - \eta(x - m e_1)|^p$ simplifies to $|\eta(x - k e_1)|^p$ on one ball and $|\eta(x - m e_1)|^p$ on the other.

    $$ ||u_k - u_m||_{L^p}^p = \int_{B(k e_1, 1)} |\eta(x - k e_1)|^p dx + \int_{B(m e_1, 1)} |\eta(x - m e_1)|^p dx = ||u_k||_{L^p}^p + ||u_m||_{L^p}^p = 2 C_1 $$

    So, for any $k \neq m$ (with $|k-m|>2$), the distance between $u_k$ and $u_m$ is constant:
    $$ ||u_k - u_m||_{L^p} = (2C_1)^{1/p} > 0 $$

    This distance does not go to zero as $k, m \to \infty$. Therefore, the sequence is not a Cauchy sequence in $L^p(\mathbb{R}^n)$. Since every convergent sequence must be Cauchy, our sequence (and any of its subsequences) cannot converge in $L^p(\mathbb{R}^n)$.
\end{enumerate}

This example perfectly illustrates the `loss of compactness' due to the unboundedness of the domain $\mathbb{R}^n$.
	
	\newpage
	\item For $u,v \in H^1(U)$, prove the integration by parts formula
	\begin{align*}
		\int_U uv_{x_i}\,dx &= \int_{\partial U} (\TRR u)(\TRR v) \nu_i\, dS- \int_U vu_{x_i}\, dx
	\end{align*}where $\TRR$ is the trace operator and $\nu_i$ is the $i\ITH$ component of the outward unit normal to $\partial U$. \\
	
	This question is a direct and natural continuation of our entire conversation. We started by discussing inequalities (Sobolev, Morrey) that define the properties of Sobolev spaces like $H^1(U) = W^{1,2}(U)$. We then discussed the concept of a \DEFINE{trace operator}, $\text{Tr}$, which gives meaning to the value of a Sobolev function on the boundary $\partial U$. We even proved that such an operator `doesn't exist' for the larger space $L^p(U)$, highlighting the importance of the derivative control provided by Sobolev spaces. This question asks us to prove the fundamental \DEFINE{integration by parts formula} for Sobolev functions, which is the primary tool that makes the trace operator so useful and is the foundation for proving many of the inequalities and PDE results we've discussed.

\textbf{Proof of Integration by Parts for $H^1(U)$ Functions}

The formula to prove is:
\[
\int_U u v_{x_i}\,dx = \int_{\partial U} (\text{Tr } u)(\text{Tr } v) \nu_i\, dS - \int_U v u_{x_i}\, dx
\]
where $u, v \in H^1(U)$ and $\nu_i$ is the $i\ITH$ component of the outward unit normal vector.

This can be rewritten by moving the second term on the right to the left side:
\[
\int_U (u v_{x_i} + v u_{x_i})\,dx = \int_{\partial U} (\text{Tr } u)(\text{Tr } v) \nu_i\, dS
\]
The integrand on the left is simply the derivative of the product, $(uv)_{x_i}$. So we need to prove:
\[
\int_U (uv)_{x_i}\,dx = \int_{\partial U} (\text{Tr } u)(\text{Tr } v) \nu_i\, dS
\]
This is the generalized version of the Divergence Theorem for Sobolev functions.

\begin{description}
\item \textbf{Step 1: The Formula for Smooth Functions}

First, recall the classical Divergence Theorem for a smooth vector field $F$ on a domain $U$ with a sufficiently regular boundary:
\[
\int_U \text{div}(F) \,dx = \int_{\partial U} F \cdot \nu \,dS
\]
If we let $F$ be the vector field that is zero in all components except for the $i$-th component, where $F_i = w$ for some smooth function $w$, then $\text{div}(F) = w_{x_i}$. The formula becomes:
\[
\int_U w_{x_i}\,dx = \int_{\partial U} w \nu_i\, dS
\]
This holds if $w$ is a smooth function, for instance, if $w \in C^1(\bar{U})$.

\item \textbf{Step 2: Density Argument}

We cannot directly apply this formula to $uv$ because if $u, v \in H^1(U)$, their product $uv$ is not necessarily in $H^1(U)$ (it is, but this requires proof). A more direct approach is to use the density of smooth functions within Sobolev spaces.

A fundamental property of Sobolev spaces is that for a domain $U$ with a reasonably nice (e.g., Lipschitz) boundary, the space of smooth functions $C^\infty(\bar{U})$ is dense in $H^1(U)$. This means that for our given $u, v \in H^1(U)$, we can find sequences of smooth functions, $\{u_k\}$ and $\{v_k\}$ in $C^\infty(\bar{U})$, such that:

\begin{itemize}
\item   $u_k \to u$ in $H^1(U)$ as $k \to \infty$
\item   $v_k \to v$ in $H^1(U)$ as $k \to \infty$
\end{itemize}

Since $u_k$ and $v_k$ are smooth, their product $u_k v_k$ is also smooth, and we can apply the classical formula from Step 1 to the function $w = u_k v_k$:
\[
\int_U (u_k v_k)_{x_i}\,dx = \int_{\partial U} (u_k v_k) \nu_i\, dS \quad (*)
\]
Our goal is to take the limit of this equation as $k \to \infty$.

\item \textbf{Step 3: Taking the Limit}

Let's analyze the limit of each side of equation $(*)$.

\begin{description}
\item \textbf{Left-Hand Side (LHS):}
    We need to show that $\int_U (u_k v_k)_{x_i}\,dx \to \int_U (uv)_{x_i}\,dx$. This is equivalent to showing $(u_k v_k)_{x_i} \to (uv)_{x_i}$ in $L^1(U)$.
    $(u_k v_k)_{x_i} = (u_k)_{x_i} v_k + u_k (v_k)_{x_i}$.
    Since $u_k \to u$ in $H^1$, we know $(u_k)_{x_i} \to u_{x_i}$ in $L^2$ and $u_k \to u$ in $L^2$. The same holds for $v_k$. Using the properties of $L^p$ spaces (specifically that convergence in $L^2$ implies convergence of products in $L^1$, a result proven with the Hölder/Cauchy-Schwarz inequality), we get:
\begin{itemize}
\item   $(u_k)_{x_i} v_k \to u_{x_i} v$ in $L^1(U)$
\item   $u_k (v_k)_{x_i} \to u v_{x_i}$ in $L^1(U)$
\end{itemize}
    Therefore, the sum converges: $(u_k v_k)_{x_i} \to (uv)_{x_i}$ in $L^1(U)$, and the integral converges.

\item \textbf{Right-Hand Side (RHS):}
    The trace operator $\text{Tr}: H^1(U) \to L^2(\partial U)$ is a bounded (and thus continuous) linear operator. This is the key property we discussed previously. Continuity means that if $u_k \to u$ in $H^1(U)$, then $\text{Tr}(u_k) \to \text{Tr}(u)$ in $L^2(\partial U)$.
    For smooth functions, $\text{Tr}(u_k) = u_k|_{\partial U}$. So we have:
\begin{itemize}
\item   $u_k|_{\partial U} \to \text{Tr}(u)$ in $L^2(\partial U)$
\item   $v_k|_{\partial U} \to \text{Tr}(v)$ in $L^2(\partial U)$
    Again, using the Hölder/Cauchy-Schwarz inequality on the boundary space $L^2(\partial U)$, the product converges in $L^1(\partial U)$:
\item   $u_k v_k|_{\partial U} \to (\text{Tr } u)(\text{Tr } v)$ in $L^1(\partial U)$
    This implies the convergence of the boundary integral:
    \[
    \int_{\partial U} u_k v_k \nu_i\, dS \to \int_{\partial U} (\text{Tr } u)(\text{Tr } v) \nu_i\, dS
    \]
\end{itemize}
\end{description}

\item \textbf{Step 4: Conclusion}

By taking the limit as $k \to \infty$ on both sides of equation $(*)$, we have shown that the LHS converges to $\int_U (uv)_{x_i}\,dx$ and the RHS converges to $\int_{\partial U} (\text{Tr } u)(\text{Tr } v) \nu_i\, dS$. Therefore, the equality must hold for the limit functions $u$ and $v$:
\[
\int_U (uv)_{x_i}\,dx = \int_{\partial U} (\text{Tr } u)(\text{Tr } v) \nu_i\, dS
\]
Expanding the left side gives the desired integration by parts formula.
\end{description}
	
	\newpage
	\item Integrate by parts to prove the interpolation inequality
	\begin{align*}
		\|Du\|_{L^2(U)}^2 \le C\|u\|_{L^2(U)}\|D^2u\|_{L^2(U)}
	\end{align*}for all $u \in H^2(U)\cap H_0^1(U)$.  (Hint: Take sequences $(v_j) \in C_c^\infty(U)$ converging to $u$ in $H_0^1(U)$ and $(w_j) \subset C^\infty(\CONJ{U})$ convergint to $u$ in $H^2(U)$.)\\
	
	This question is an excellent synthesis of our entire conversation. It asks to prove an \DEFINE{interpolation inequality}, a concept we've already explored, which falls under the broad category of inequalities like \DEFINE{Cauchy-Schwarz}, \DEFINE{Sobolev}, and \DEFINE{Morrey} that define the structure of function spaces. The proof requires using \DEFINE{integration by parts for Sobolev functions}, a formula we just derived. Furthermore, the function space in question, $H^2(U) \cap H_0^1(U)$, combines the concepts of higher-order \DEFINE{regularity} ($H^2$) and zero boundary conditions ($H_0^1$), both of which are central themes we have discussed, particularly in relation to the \DEFINE{trace theorem}.

Let's prove the inequality for $u \in H^2(U) \cap H_0^1(U)$.

\textbf{Proof}`

Our goal is to prove that there exists a constant $C$ such that:
\[
\|Du\|_{L^2(U)}^2 \le C \|u\|_{L^2(U)}\|D^2u\|_{L^2(U)}
\]

\begin{description}
\item \textbf{Step 1: Start with the integral form}

The left-hand side is the integral of $|Du|^2$. We can write this using summation notation as:
\[
\|Du\|_{L^2}^2 = \int_U |Du|^2 \,dx = \int_U \sum_{i=1}^n u_{x_i} u_{x_i} \,dx
\]
where $u_{x_i} = \frac{\partial u}{\partial x_i}$.

\item \textbf{Step 2: Apply Integration by Parts}

As we proved in our previous discussion, the integration by parts formula holds for functions in $H^1(U)$. Since $u \in H^2(U)$, both $u$ and its derivatives $u_{x_i}$ are in $H^1(U)$. We can apply the formula to the product of $u$ and $u_{x_i}$:
\[
\int_U u_{x_i} u_{x_i} \,dx = \int_U (u_{x_i}) u_{x_i} \,dx = \int_{\partial U} (\text{Tr } u_{x_i})(\text{Tr } u) \nu_i \,dS - \int_U u \, u_{x_i x_i} \,dx
\]
Here, $\text{Tr}$ is the trace operator and $\nu_i$ is the $i$-th component of the outward normal vector.

\item \textbf{Step 3: Use the Boundary Condition}

The condition $u \in H_0^1(U)$ is crucial. By definition, this means that the trace of $u$ on the boundary is zero, i.e., $\text{Tr } u = 0$. Therefore, the boundary integral vanishes:
\[
\int_{\partial U} (\text{Tr } u_{x_i})(\text{Tr } u) \nu_i \,dS = 0
\]
This simplifies our equation significantly:
\[
\int_U u_{x_i} u_{x_i} \,dx = - \int_U u \, u_{x_i x_i} \,dx
\]

\item \textbf{Step 4: Sum over all derivatives and apply Cauchy-Schwarz}

Now, we sum this result over all $i = 1, \dots, n$:
\[
\|Du\|_{L^2}^2 = \sum_{i=1}^n \int_U u_{x_i}^2 \,dx = -\sum_{i=1}^n \int_U u \, u_{x_i x_i} \,dx = -\int_U u \left( \sum_{i=1}^n u_{x_i x_i} \right) \,dx = -\int_U u (\Delta u) \,dx
\]
where $\Delta u$ is the Laplacian of $u$. Now, we take the absolute value and apply the \DEFINE{Cauchy-Schwarz inequality}, which we identified at the beginning of our conversation as the fundamental tool for $L^2$ spaces:
\[
\|Du\|_{L^2}^2 = \left| -\int_U u (\Delta u) \,dx \right| \le \int_U |u| |\Delta u| \,dx \le \|u\|_{L^2} \|\Delta u\|_{L^2}
\]

\item \textbf{Step 5: Relate the Laplacian to the full second derivative}

The norm of the Laplacian, $\|\Delta u\|_{L^2}$, can be bounded by the norm of the full Hessian matrix of second derivatives, $\|D^2u\|_{L^2}$. The norm $\|D^2u\|_{L^2}$ is defined as $\left( \sum_{i,j=1}^n \|u_{x_i x_j}\|_{L^2}^2 \right)^{1/2}$. By the triangle inequality:
\[
\|\Delta u\|_{L^2} = \left\| \sum_{i=1}^n u_{x_i x_i} \right\|_{L^2} \le \sum_{i=1}^n \|u_{x_i x_i}\|_{L^2} \le \sqrt{n} \left( \sum_{i=1}^n \|u_{x_i x_i}\|_{L^2}^2 \right)^{1/2} \le \sqrt{n} \|D^2u\|_{L^2}
\]
Substituting this into our inequality, we get:
\[
\|Du\|_{L^2}^2 \le \|u\|_{L^2} (\sqrt{n} \|D^2u\|_{L^2})
\]
Setting $C = \sqrt{n}$ (a constant that depends only on the dimension), we arrive at the final result:
\[
\|Du\|_{L^2(U)}^2 \le C \|u\|_{L^2(U)} \|D^2u\|_{L^2(U)}
\]
\end{description}
\end{enumerate}

\section{Inequalities}
\begin{description}
	\item \textbf{Gagliardo-Nirenberg-Sobolev}
	
	Assume $1\le p < n$.  Ther exists a constant $C$, depending onlyon $p$n and $n$ , such that 
	\begin{align*}
		\|u\|_{L^{p^*}(\R^n)} \le C\|Du\|_{L^p(\R^n)}
	\end{align*} for all $u \in C_c^1(\R^n)$.
	
	\item \textbf{Theorem 2 (EStimates for $W^{1,p}, 1 \le p < n)$}
	
	Let $U$ b3 bounded, open subset of $\R^n$, and suppose $\partial U$ is $C^1$.  Assume $1 \le p < n$, and $u\in W^{1,p}(U)$.  Then $u \in L^{p^*}(U)$, whtih the estimate 
	\begin{align*}
		\|u\|_{L^{p^*}(U)},
	\end{align*}the constant $C$ depending only $p,n$, and $U$.
	
	\item \textbf{Morrey's Inequality}
	
	Assume $n < p \le \infty$.  Then ther exists a constant $C$, depending only on $p$ and $n$, such that 
	\begin{align*}
		\|u\|_{C^{0,\gamma}(\R^n)} \le C\|u\|_{W^{1,p}(\R^n)}
	\end{align*}for all $u \in C^1(\R^n)$, where
	\begin{align*}
		\gamma := 1-n/p.
	\end{align*}
	
	\item \textbf{Theorem 5: Estimates for $W^{1,p}, n < p \le \infty$.}
	
	Let $U$ be a bounded, open subset of $\R^n$, and suppose $\partial U$ is $C^1$.'  Assume $n < p \le \infty$ and $u \in W^{1,p}(U)$.  Then $u$ has a version $u^*\in C^{0,\gamma}(\CONJ{U})$, for $\gamma = 1-\frac{n}{p}$, with the estimate 
	\begin{align*}
		\|u^*\|_{C^{0,\gamma}(\CONJ{U})} \le C\|u\|_{W^{1,p}(U)}.
	\end{align*}The constant $C$ depends only on $p,n$, and $U$.
	
	\item \textbf{Theorem 6: General Sobolev Inequalities).}
	
	Let $U$ be a bounded open subset on $\R^n$, with a $C^1$ boundary.  Assume $u \in W^{k,p}(U)$.
	\begin{enumerate}[label=(\roman*)]
		\item If 
		\begin{align*}
			k < \frac{n}{p}
\end{align*}	 then $u \in L^q(U)$, where
\begin{align*}
	\frac{1}{q}= \frac{1}{p}-\frac{k}{n}
\end{align*}We have addition the estimate
\begin{align*}
	\|u\|_{L^q(U)} \le C \|u\|_{W^{k,p}(U)},
\end{align*}the constant $C$ depending only on $k,p,n$ and $U$.
		\item If
		\begin{align*}
			k > \frac{n}{p}
		\end{align*}then $u \in C^{k-\frac{n}{p}-1,\gamma}(\CONJ{U})$, where
		\begin{align*}
			\gamma &= \BINDEF{\frac{n}{p}+1-\frac{n}{p}, & \text{if $\frac{n}{p}$ is not an integer}}{\text{any positive number } < 1, &\text{if $\frac{n}{p}$ is an integer}}
		\end{align*}We have an addition the estimate
depending only on $k,p,n,\gamma$ and $U$.
		\begin{align*}
			\gamma &= \BINDEF{\frac{n}{p}+1-\frac{n}{p}, & \text{if $\frac{n}{p}$ is not an integer}}{\text{any positive number } < 1, &\text{if $\frac{n}{p}$ is an integer}}
		\end{align*}We have an addition the estimate
		\begin{align*}
			\|u\|_{C^{k-\frac{n}{p}-1,\gamma}(\CONJ{U})} \le C\|u\|_{W^{k,p}(U)},
		\end{align*}the constant $C$ depending only on $k,p,n,\gamma$ and $U$.
	\end{enumerate}
	
\end{description}
\end{document}
